  

%%%%%%%%%%%%%%Empieza la presentacion%%%%%%%%%%%%%%%%%%%%%%%%%%%%%%%%%%%%

%\begin{frame}
%\frametitle{Esquema}
 %   \tableofcontents
%\end{frame}

\section{Objetivos}

\begin{frame}
\frametitle{Objetivos del tema}

\begin{itemize}
\item Aprender a utilizar técnicas  de diseño de filtros generales (paso baja, paso alta, paso banda y supresor de banda):

\begin{itemize}
\normalsize
\visible<2->{
\item Continuos: 
\begin{itemize}
\small
\item Obtener únicamente la función de transferencia del filtro (no vamos a diseñar el sistema analógico que lo representa)
\item Aproximaciones a los filtros ideales
\end{itemize}
}

\visible<3->{
\item Discretos:
\begin{itemize}
\small
\item Filtros IIR: a partir de las aproximaciones de filtros continuos
\item Filtros FIR: a partir de la transformada de Fourier con ventana
\end{itemize}

}

\end{itemize}

\visible<4->{\item Aprender a diseñar filtros específicos/a medida}

\end{itemize}

\end{frame}

